\begin{textsample}{POS Dim 4 – 2016\_c6 – Score 5.00 – 2016\_c6/enade\_2016\_tec\_agro\_1c1f311707.txt}  \label{ex:f4_pos_054}
CAPÍTULO 3 – ANÁLISE TÉCNICA DA PROVA Este capítulo tem por objetivo apresentar o desempenho dos estudantes concluintes de Tecnologia em Agronegócio no Enade/2016.
Para isso, foram calculadas as estatísticas básicas da prova em seu todo ( seção 3.1.1 ), bem como as estatísticas relacionadas aos Componentes de Formação Geral ( seção 3.1.2 ) e de Conhecimento Específico da Área ( seção 3.1.3 ).
Dadas as suas características, foram analisadas, em separado, as questões objetivas ( seção 3.2 ) e as questões discursivas ( seção 3.3 ).
Para as questões objetivas, foram disponibilizados os índices de facilidade e de discriminação ponto-bisserial, também em separado, para os Componentes de Formação Geral ( seção 3.2.1 ) e de Conhecimento Específico ( seção 3.2.2 ).
De cada componente, uma das questões foi escolhida para exemplificar a análise gráfica, relacionando as alternativas escolhidas pelos estudantes ( inclusive o gabarito ) com o número de acertos no componente.
O Anexo I apresenta a íntegra da análise gráfica para todas as questões objetivas.
Para cada uma das questões discursivas, os conteúdos dos tipos mais comuns de respostas dos estudantes são apresentados e comparados com o padrão de respostas esperado ( ver Anexo VIII com o padrão de respostas ).
Tomando como base as duas questões discursivas do Componente de Formação Geral, a seção 3.3.1.6 apresenta comentários sobre a correção das respostas com respeito à Língua Portuguesa.
Nas Tabelas 3.1 e 3.3 são apresentados o tamanho da população inscrita, ausente, presente e percentual de faltosos, respectivamente, por Grande Região e por Categoria Administrativa, Organização Acadêmica da IES e Modalidade de Ensino.
As demais tabelas apresentam as seguintes estatísticas das notas15 : média do desempenho na prova, erro padrão da média, desvio padrão, nota mínima, mediana e nota máxima para cada um de seus componentes.
Tais estatísticas contemplam o total de estudantes concluintes da Área de Tecnologia em Agronegócio inscritos e presentes à prova do Enade/2016, tendo em vista agregações, ou por Grandes Regiões e o país como um todo, ou por Categoria Administrativa, Organização Acadêmica da IES e Modalidade de Ensino.
Em relação aos gráficos de distribuição de notas, o intervalo considerado foi de 10 unidades, aberto à esquerda e fechado à direita, com exceção do primeiro intervalo, [ 0 ; 10 ], fechado em ambos os extremos.
Para os gráficos de distribuição das notas das 15 Uma definição dessas estatísticas pode ser encontrada no Glossário. 35 3.1 - CONDIÇÃO DE PRESENÇA E ESTATÍSTICAS BÁSICAS DA PROVA Esta seção apresenta a condição de presença, além das estatísticas selecionadas e histogramas da nota geral ( 3.1.1 ) e de cada componente : Formação Geral ( 3.1.2 ) e Conhecimento Específico ( 3.1.3 ).
São também apresentadas estatísticas selecionadas de subpopulações caracterizadas por Grande Região, Categoria Administrativa, Organização Acadêmica e Modalidade de Ensino. 3.1.1 - Condição de Presença e Estatísticas Básicas Gerais A Tabela 3.1 apresenta a Condição de Presença na prova, por Grande Região, dos estudantes concluintes de Tecnologia em Agronegócio.
A população total de Inscritos foi de 1.701.
Destes, 1.419 estiveram Presentes, sendo 16,6 \% o índice de não comparecimento.
A Região de maior ausência foi a Norte ( 18,9 \%, ) e a de menor ausência foi a Sul ( 11,4 \% ).
Já a Tabela 3.2 apresenta as Estatísticas Básicas da prova, por Grande Região, dos estudantes concluintes de Tecnologia em Agronegócio.
A Média das notas da prova como um todo ( nas seções seguintes serão analisados os Componentes de Formação Geral e de Conhecimento Específico ) foi 42,4, sendo que os alunos da região Norte obtiveram a Média mais baixa ( 38,4 ), e os da região Sul obtiveram a Média mais alta ( 44,9 ).
As demais médias foram : 39,6 na região Nordeste, 40,7 na região Centro-Oeste e 43,0 na região Sudeste.
O Desvio padrão para o Brasil como um todo foi 12,5, sendo o maior Desvio padrão encontrado na região Sul ( 12,7 ) e o menor, na região Sudeste ( 12,0 ), indicando uma dispersão um pouco menor das notas desta última região.
A região que obteve a maior nota Máxima foi a Sul ( 81,3 ), ao passo que a região que atingiu a menor nota Máxima foi a Norte ( 66,0 ).
A Mediana do Brasil como um todo foi 42,4, sendo a maior Mediana obtida na região Sul ( 45,6 ), e a menor obtida na Nordeste ( 38,8 ).
A menor nota Mínima foi obtida na região Nordeste ( 3,4 ), e a maior, na região Sudeste ( 14,8 ).
Considerando -se as notas segundo Grande Região, observa -se que existe diferença estatisticamente significativa ao nível de 95 \% entre a menor Média, obtida na região Norte ( 38,4 ), e a maior Média, obtida na região Sul (44,9)17.
O comportamento das notas dos estudantes de todo o Brasil pode ser observado no Gráfico 3.1 que apresenta um histograma com a distribuição das mesmas.
Essa é uma distribuição unimodal com a moda no intervalo ( 40 ; 50 ], seguido pelo intervalo ( 30, 40 ]. 17 Todas as comparações deste capítulo utilizam os intervalos de 95 \%.
Os erros-padrão da média que possibilitam os testes estão disponíveis nas tabelas.
A Tabela 3.3 apresenta informações referentes à Condição de Presença desagregadas por Categoria Administrativa, por Organização Acadêmica e por Modalidade de Ensino.
Da população total de Inscritos, 818 são de IES Privadas e 883 de IES Públicas.
Em relação à Organização Acadêmica, a maior participação foi obtida por estudantes de Faculdades ( 871 ), seguido de Universidades ( 625 ) e Centros Universitários ( 205 ).
A grande maioria dos concluintes inscritos foi de Educação Presencial, 1.471, ao passo que a modalidade a distância somou 230.
A Categoria Administrativa de maior ausência foi a Pública ( 19,8 \% ), e entre as Organizações Acadêmicas foi a das Universidades ( 18,7 \% ), que ficou também acima da taxa de ausência nacional de 16,6 \%.
Na Modalidade de Educação a Distância, a ausência foi de 14,8 \%, e na Educação Presencial, 16,9 \%.
Já a Tabela 3.4 apresenta informações referentes à Média da nota final desagregadas por Categoria Administrativa, por Organização Acadêmica e por Modalidade de Ensino.
A Média das notas da prova como um todo foi 42,4.
Em relação à Categoria Administrativa, os estudantes das IES Privadas obtiveram Média um pouco mais alta ( 42,5 ) que a do Brasil como um todo, e os das IES Públicas obtiveram Média ligeiramente mais baixa ( 42,3 ) que a Média nacional.
Observa -se que não existe diferença estatisticamente significativa entre as médias das notas das IES Públicas e Privadas.
A diferença entre as médias das regiões Sul e Norte ( 6,5 ), a maior e a menor Média, é superior à diferença entre IES Públicas e Privadas ( 0,2 ), caracterizando uma menor diversidade administrativa do que regional.
O Desvio padrão para as IES Públicas ( 12,2 ) foi inferior ao do Brasil como um todo ( 12,5 ), indicando uma dispersão um pouco menor das notas nesta Categoria Administrativa.
No tocante à Organização Acadêmica, os Centros Universitários obtiveram Média mais alta que a nacional ( 45,2 e 42,4, respectivamente ).
A média das Faculdades foi 42,1 e a das Universidades foi 41,9.
Constata -se que existe diferença estatisticamente significativa ao nível de 95 \% entre a maior Média ( Centros Universitários ) e a menor Média ( Universidades ), no que diz respeito às Organizações Acadêmicas.
A Média da Modalidade de Ensino de Educação a Distância ( 47,0 ) foi superior à Média da Educação Presencial ( 41,6 ).
Também houve diferença estatisticamente significativa ao nível de 95 \% entre essas duas médias. 3.1.2 - Estatísticas Básicas no Componente de Formação Geral A Tabela 3.5 apresenta as Estatísticas Básicas em relação ao componente da prova que avalia a Formação Geral dos estudantes concluintes.
Os alunos de todo o Brasil obtiveram desempenho médio de 40,9.
Quanto à variabilidade, o Desvio padrão das notas dos estudantes do Brasil, como um todo, foi 15,1.
A maior Média foi obtida na região Sudeste ( 42,0 ), e a menor, na região Norte ( 37,8 ).
As demais médias foram : 40,0 na região Sul, 40,4 na região Centro-Oeste e 41,8 na região Nordeste.
Já o maior Desvio padrão foi obtido na região Sul ( 15,4 ), e o menor, na região Nordeste ( 14,1 ).
Os demais desvios padrões foram : 14,9 na região Norte, 15,1 na região Centro-Oeste e 15,2 da região Sudeste.
A maior nota no Componente de Formação Geral da prova do Enade/2016 foi 91,2, obtida por, pelo menos, um aluno na região Centro-Oeste.
A menor nota Máxima foi obtida na região Norte ( 71,2 ).
A Mediana do Brasil, como um todo, foi 40,4, sendo a menor Mediana encontrada na região Norte ( 37,8 ), e a maior, encontrada na região Nordeste ( 42,6 ).
A nota Mínima nesta parte foi zero em todas as regiões, sem exceção.
Considerando -se as notas segundo Grande Região, observa -se que não existe diferença estatisticamente significativa entre qualquer Média das notas do Componente de Formação Geral.
O Gráfico 3.2 propicia a avaliação do desempenho dos estudantes do Componente de Formação Geral, a partir do histograma da distribuição das notas correspondentes.
A distribuição é unimodal, com moda em ( 30 ; 40 ], enquanto na prova, como um todo ( Gráfico 3.1 ), a moda foi alcançada no intervalo ( 40 ; 50 ].
Nota -se, ainda que, no Gráfico 3.2, as notas apresentam maior dispersão do que no Gráfico 3.1 ( distribuição das notas da prova ), confirmada pela comparação dos desvios padrões : 12,5 para a nota da prova como um todo e 15,1 para o Componente de Formação Geral.
Na Tabela 3.6 são apresentadas as informações referentes ao desempenho dos concluintes do Componente de Formação Geral, em diferentes agregações : Categoria Administrativa, Organização Acadêmica e Modalidade de Ensino.
Observa -se que existe diferença estatisticamente significativa entre as médias dos tipos de Categoria Administrativa.
A maior Média foi obtida por estudantes de IES Públicas ( 42,8 ), com uma diferença estatisticamente significativa da obtida por estudantes de IES Privadas ( 38,9 ).
Nota -se que também não há diferença estatisticamente significativa entre as médias dos três tipos de Organização Acadêmica, sendo o maior valor observado nas Faculdades e o menor nos Centro Universitários.
No que se refere à Modalidade de Ensino, a Educação Presencial apresentou maior média do que a Educação a Distância ( médias 41,1 e 39,4, respectivamente ).
Tais médias também não apresentam diferenças estatisticamente significativas. 3.1.3 - Estatísticas Básicas do Componente de Conhecimento Específico A Tabela 3.7 apresenta as Estatísticas Básicas referentes ao Componente de Conhecimento Específico da área de Tecnologia em Agronegócio.
A Média do desempenho dos alunos do Brasil, como um todo, foi 42,9.
A maior Média foi obtida na região Sul ( 46,6 ), e a menor, na região Norte ( 38,6 ).
As demais médias foram : 38,9 na região Nordeste, 40,7 na região Centro-Oeste e 43,4 na região Sudeste.
Quanto à variabilidade das notas, o Desvio padrão do Brasil, como um todo, foi 14,1, sendo o maior Desvio padrão observado na região Nordeste ( 14,9 ), e o menor, na região Sudeste ( 13,4 ).
Os demais desvios foram : 13,7 na região Norte, 13,8 na região Centro-Oeste e 14,3 na região Sul.
A Mediana das notas dos estudantes de todo o Brasil foi 43,0.
A maior Mediana ocorreu na região Sul ( 47,4 ), e a menor, na região Nordeste ( 36,3 ).
As demais medianas foram : 39,5 na região Norte, 41,5 na região Centro-Oeste e 43,0 na região Sudeste.
A nota Máxima do Brasil, como um todo, foi 89,0, sendo obtida por, pelo menos, um aluno na região Sul.
As demais notas máximas foram : 70,8 na região Nordeste, 76,5 na região Norte, 78,0 na região Centro-Oeste e 78,8 na região Sudeste.
A nota Mínima foi 2,0 e foi obtida na região Nordeste.
A região Centro-Oeste teve nota Mínima 4,0, nas regiões Sudeste e Sul a Mínima foi 10,0 e na Norte, foi 10,5.
Observa -se que existe diferença estatisticamente significativa da menor Média das notas do Componente de Conhecimento Específico da região Norte ( 38,6 ) em relação à maior Média da região Sul ( 46,6 ).
Assim como os Gráficos 3.1 e 3.2, o Gráfico 3.3, apresentado a seguir, permite uma avaliação do desempenho de concluintes de Tecnologia em Agronegócio em relação ao Componente de Conhecimento Específico com um histograma da distribuição das notas correspondentes.
Esta também é uma distribuição unimodal, e o grupo modal é o ( 40 ; 50 ], o mesmo grupo modal da prova como um todo e acima do grupo modal para a Formação Geral.
A Tabela 3.8 apresenta uma comparação dos resultados em relação à Categoria Administrativa, à Organização Acadêmica e a Modalidade de Ensino, agora levando em conta o desempenho de alunos do Componente de Conhecimento Específico da prova.
No que se refere à Organização Acadêmica, a maior Média foi dos Centros Universitários ( 47,2 ), vindo a seguir Universidades e Faculdades, ambas Organizações Acadêmicas com 42,2.
O maior Desvio padrão, e acima do valor para o Brasil, como um todo, foi o dos Centros Universitários ( 15,2 ).
Os Centros Universitários também obtiveram as maiores notas Máxima ( 89,0 ) e Mediana ( 47,0 ).
As Universidades obtiveram nota Máxima igual a 80,0 e as Faculdades obtiveram nota Máxima igual a 78,8.
As medianas foram : 42,5 nas Universidades e 42,4 nas Faculdades.
As notas Mínimas foram 10,0, 4,0 e 2,0 para os Centros Universitários, as Universidades e as Faculdades, respectivamente.
Observa -se que existe diferença estatisticamente significativa ao nível de 95 \% no Componente de Conhecimento Específico entre a maior 3.2 - ANÁLISE DAS QUESTÕES OBJETIVAS Esta seção apresenta estatísticas selecionadas e histogramas das Questões Objetivas de Formação Geral ( 3.2.1 ).
São também apresentadas e comparadas as médias das subpopulações caracterizadas por Grande Região. 3.2.1 - Componente de Formação Geral A Tabela 3.9 apresenta as Estatísticas Básicas relativas às oito questões objetivas do componente da prova que abrange a Formação Geral dos estudantes.
A Média do Brasil foi 33,5.
A menor Média foi encontrada na região Norte ( 30,8 ), e a maior, na região Sudeste ( 34,9 ).
As demais médias foram : 31,9 na região Sul, 33,2 na região Centro-Oeste e 34,4 na região Nordeste.
O Desvio padrão do Brasil foi 18,8, sendo o maior Desvio padrão encontrado na região Sul ( 19,2 ), e o menor, na região Nordeste ( 18,0 ).
Os demais desvios foram : 18,1 na região Norte, 18,3 na região Centro-Oeste e 19,0 na região Sudeste. e a menor nota ( a nota dos Centros Universitários e a nota das Faculdades, respectivamente ).
Quanto à Categoria Administrativa, observa -se que não existe diferença estatisticamente significativa entre as médias das IES Públicas ( 42,1 ) e IES Privadas ( 43,7 ).
A maior Média foi obtida por alunos de IES Privadas de ensino.
No que toca à Modalidade de Ensino, há diferença estatisticamente significativa entre as médias da Educação Presencial ( 41,8 ) e da Educação a Distância ( 49,5 ). 3.2 - ANÁLISE DAS QUESTÕES OBJETIVAS Esta seção apresenta estatísticas selecionadas e histogramas das Questões Objetivas de Formação Geral ( 3.2.1 ).
São também apresentadas e comparadas as médias das subpopulações caracterizadas por Grande Região. 3.2.1 - Componente de Formação Geral A Tabela 3.9 apresenta as Estatísticas Básicas relativas às oito questões objetivas do componente da prova que abrange a Formação Geral dos estudantes.
A Média do Brasil foi 33,5.
A menor Média foi encontrada na região Norte ( 30,8 ), e a maior, na região Sudeste ( 34,9 ).
As demais médias foram : 31,9 na região Sul, 33,2 na região Centro-Oeste e 34,4 na região Nordeste.
O Desvio padrão do Brasil foi 18,8, sendo o maior Desvio padrão encontrado na região Sul ( 19,2 ), e o menor, na região Nordeste ( 18,0 ).
Os demais desvios foram : 18,1 na região Norte, 18,3 na região Centro-Oeste e 19,0 na região Sudeste.
As medianas do Brasil como um todo e de três das cinco regiões foi 37,5 — as duas regiões com mediana diferente foram a Norte e a Sul, cuja Mediana foi 25,0 em cada uma delas.
A nota Máxima do Brasil foi 100,0 e ocorreu nas regiões Sudeste e Centro-Oeste — as outras regiões tiveram a mesma nota Máxima igual a 87,5.
A nota Mínima 0,0 ocorreu em todas as regiões.
A Tabela 3.10 apresenta o Índice de Facilidade e o Índice de Discriminação ( ponto-bisserial ) para cada uma das questões objetivas do Componente de Formação Geral.
Quanto ao Índice de Facilidade, foram usadas as seguintes cores para diferenciar o nível de dificuldade da questão :  Azul para as questões classificadas com índice Muito fácil ( \&gt ; =0,86 ), verde para as questões classificadas com índice Fácil ( 0,61 a 0,85 ), amarelo para as questões classificadas com Médio ( 0,41 a 0,60 ), vermelho para as questões classificadas com Difícil ( 0,16 a 0,40 ) e roxo para as questões classificadas com Muito difícil ( \&lt ; =0,15 ).
Já quanto ao Índice de Discriminação, foram usadas as seguintes cores para qualificar a questão :  As questões classificadas com índice Fraco receberam a cor vermelho ( \&lt ; =0,19 ), as classificadas com Médio receberam a cor amarelo ( 0,20 a 0,29 ), as classificadas com Bom receberam a cor verde ( 0,30 a 0,39 ) e as classificadas com Muito bom ( \&gt ; =0,40 ) receberam a cor azul.
As questões objetivas do Componente de Formação Geral, segundo o Índice de Facilidade, foram assim avaliadas : das oito questões, nenhuma teve o Índice de Facilidade classificado como Muito fácil e uma questão foi tida como Fácil, por ter índice de acertos 0,69.
Duas questões foram consideradas com índice de dificuldade Médio, situando -se no intervalo entre 0,41 e 0,60 do Índice de Facilidade, ou seja, houve entre 41,0 \% e 60,0 \% de acertos, enquanto quatro questões foram classificadas na categoria Difícil, situando -se no intervalo entre 0,16 e 0,40.
Por fim, uma questão apresentou menos de 15 \% de acertos ( 9 \% ), razão pela qual seria classificada como Muito difícil.
Como já comentado, para análise das questões objetivas, relativas à Formação Geral, segundo o poder de discriminação, utilizou -se o Índice de Discriminação ( ponto-bisserial ).
Nesta análise, as questões foram assim avaliadas : seis das oito questões apresentaram índice acima ou igual a 0,40 e, assim, foram classificadas com índice Muito bom para esse grupo de estudantes.
Uma questão teve Índice de Discriminação Bom, com valor entre 0,30 e 0,39, e uma questão foi classificada com índice de discriminação Médio ( entre 0,20 e 0,29 ).
Nenhuma questão teve nível Fraco de discriminação para esse grupo de estudantes.
O Índice de Facilidade variou de 0,09 a 0,69, e o de Discriminação, de 0,27 a 0,52.
As seis questões com Índice de Discriminação Muito bom figuraram entre níveis de dificuldade Fácil, Médio e Difícil : uma classificada na categoria Fácil ( a questão 6 ) do Índice de Facilidade, duas, na categoria Médio ( questões 1 e 3 ) e três na categoria Difícil ( questões 2, 4 e 7 ).
Em particular, a questão 7 foi a que apresentou o maior poder discriminatório, com índice 0,52, e foi considerada Difícil em termos de facilidade, com uma proporção de 0,32 acertos.
O máximo de acertos foi alcançado pela questão 6 com um Índice de Facilidade de 0,69.
A questão de número 8 apresentou um Índice de Facilidade de 0,19, ou seja, menos do que um quinto dos estudantes conseguiu resolvê -la, dentro do universo de participantes.
Seu Índice de Discriminação foi Bom ( 0,37 ).
Já a questão 4 obteve Índice de Discriminação Muito bom, 0,40, e seu Índice de Facilidade foi Difícil ( 0,16 ).
Difícil, situando -se no intervalo entre 0,16 e 0,40.
Por fim, uma questão apresentou menos de 15 \% de acertos ( 9 \% ), razão pela qual seria classificada como Muito difícil.
Como já comentado, para análise das questões objetivas, relativas à Formação Geral, segundo o poder de discriminação, utilizou -se o Índice de Discriminação ( ponto-bisserial ).
Nesta análise, as questões foram assim avaliadas : seis das oito questões apresentaram índice acima ou igual a 0,40 e, assim, foram classificadas com índice Muito bom para esse grupo de estudantes.
Uma questão teve Índice de Discriminação Bom, com valor entre 0,30 e 0,39, e uma questão foi classificada com índice de discriminação Médio ( entre 0,20 e 0,29 ).
Nenhuma questão teve nível Fraco de discriminação para esse grupo de estudantes.
O Índice de Facilidade variou de 0,09 a 0,69, e o de Discriminação, de 0,27 a 0,52.
As seis questões com Índice de Discriminação Muito bom figuraram entre níveis de dificuldade Fácil, Médio e Difícil : uma classificada na categoria Fácil ( a questão 6 ) do Índice de Facilidade, duas, na categoria Médio ( questões 1 e 3 ) e três na categoria Difícil ( questões 2, 4 e 7 ).
Em particular, a questão 7 foi a que apresentou o maior poder discriminatório, com índice 0,52, e foi considerada Difícil em termos de facilidade, com uma proporção de 0,32 acertos.
O máximo de acertos foi alcançado pela questão 6 com um Índice de Facilidade de 0,69.
A questão de número 8 apresentou um Índice de Facilidade de 0,19, ou seja, menos do que um quinto dos estudantes conseguiu resolvê -la, dentro do universo de participantes.
Seu Índice de Discriminação foi Bom ( 0,37 ).
Já a questão 4 obteve Índice de Discriminação Muito bom, 0,40, e seu Índice de Facilidade foi Difícil ( 0,16 ).
O Gráfico 3.4, para exemplificar, analisa o comportamento da questão de número 2 de Formação Geral.
Trata -se de uma questão considerada Difícil, com relação à facilidade e a que obteve o terceiro maior Índice de Discriminação dessa parte da prova.
Neste gráfico, cada uma das cinco curvas representa o percentual de respostas em determinada alternativa da questão, em função do número de acertos dos estudantes nessa parte da prova ( Formação Geral/Múltipla Escolha ), antes de possíveis eliminações pelo critério do ponto-bisserial.
A curva em azul corresponde à alternativa A, a alternativa correta para esta questão.
Observa -se que entre os estudantes com menor número de acertos, nessa parte do exame, a situação mais frequente foi a escolha da alternativa E ( em vermelho ), incorreta.
Por exemplo, entre os estudantes que acertaram três questões, 33,9 \% escolheram a alternativa E, 32,8 \% escolheram a alternativa B ( em verde ), 1,4 \% escolheu a alternativa C ( preto ) e 1,7 \% a D ( roxo ).
Entre os que acertaram três das questões de múltipla escolha de formação geral, 0,0 \% deixou a questão em branco e 0,0 \% marcou mais de uma alternativa, o que invalidaria a questão.
À medida em que o número de acertos aumenta, indicando desempenho melhor nessa parte da prova, aumenta concomitantemente a proporção de estudantes que selecionaram a alternativa correta A, atingindo 100 \% para os estudantes com 8 acertos.
Essa análise permite verificar como a questão discriminou os grupos de desempenho, justificando o alto índice obtido na questão.
Cumpre notar que não é possível inferir deste gráfico nem o índice de facilidade ( que seria uma média da proporção ponderada pela quantidade de alunos com cada uma das notas ), nem o índice de discriminação ponto-bisserial, por razão equivalente.
No caso extremo no qual a grande concentração dos acertos dos alunos fosse abaixo de 4, o índice de facilidade seria obrigatoriamente abaixo de 40 \% ( neste exemplo ).
Caso a concentração fosse em 6 acertos ou mais, o índice seria obrigatoriamente acima de 60 \%.
Os gráficos relativos às demais questões de Formação Geral constam do Anexo I. 3.2.2 - Componente de Conhecimento Específico A Tabela 3.11 apresenta as Estatísticas Básicas em relação às questões objetivas do Componente de Conhecimento Específico da prova de Tecnologia em Agronegócio por Grande Região.
A Média do Brasil deste componente foi de 42,6.
A menor Média foi observada na região Norte ( 38,3 ), e a maior, na região Sul ( 46,3 ).
O Desvio padrão de todo o Brasil foi 15,4, sendo o menor Desvio padrão encontrado nas regiões Norte e Sudeste ( 14,8 ), e o maior, na região Nordeste ( 16,4 ).
A Mediana de todo o Brasil foi 41,2, a mesma encontrada em duas das regiões : Sudeste e Centro-Oeste.
Nas regiões Norte e Nordeste a Mediana foi menor ( 35,3 em cada uma delas ) e na região Sul, maior ( 47,1 ).
A nota Máxima da prova foi 88,2, obtida nas questões objetivas do Componente de Conhecimento Específico, por, pelo menos, um aluno na região Sul.
Nas demais regiões, a nota Máxima da prova foi : 70,6 na região Nordeste, 76,5 na região Centro-Oeste e 82,4 nas regiões Norte e Sudeste.
A nota Mínima foi zero nas regiões Nordeste e Centro-Oeste, e 5,9 nas demais regiões.
A Tabela 3.12 apresenta os Índices de Facilidade e Discriminação ( ponto-bisserial ) das questões objetivas do Componente de Conhecimento Específico para os estudantes de Tecnologia em Agronegócio.
Para facilitar a diferenciação das questões, usaremos as mesmas cores da Tabela 3.10 para as diferentes classificações dos Índices de Facilidade e de Discriminação.
A partir do Índice de Facilidade obtido, pode -se concluir que quase três quartos das questões objetivas da prova foram considerados, pelo menos, Difícil : das 27 questões válidas, 20 foram classificadas como Difícil ou como Muito difícil.
Duas questões foram classificadas como Fácil, e outras cinco consideradas como Médio.
Já quanto aos índices de discriminação das questões objetivas do Componente de Conhecimento Específico da prova, tem -se como resultado a seguinte classificação : seis das 27 questões foram consideradas como boas, enquanto uma delas teve Índice de Discriminação Muito bom.
Assim, para um pouco mais de um quarto das questões – sete em 27 – os Índices de Discriminação foram Bom ou Muito bom.
Dentre as demais, 10 delas foram classificadas como Médio e outras 10 como Fraco, sendo 20, por conseguinte, a quantidade de questões nos dois patamares mais baixos de discriminação.
Constata -se, assim, que a prova – no que se refere ao Componente de Conhecimento Específico – possuía pouca capacidade de discriminar entre aqueles que dominam ou não o conteúdo.
Dentre as questões que alcançaram os maiores índices de discriminação, uma delas, a de número 32, foi classificada com Índice de Discriminação Muito bom, situando -se no intervalo de 0,40 a 0,50 do índice, e foi classificada na categoria Médio, quanto ao Índice de Facilidade.
A questão de número 30 foi a mais difícil dentre as 27 questões específicas, com baixo Índice de Facilidade, apenas 8 \% de acertos.
Essa questão apresentou poder discriminatório igualmente baixo, – 0,02, o que comprova ter sido esta questão a mais difícil para os estudantes.
Destacam -se, também, as questões 19 e 21, com Índices de Facilidade 0,11 e 012, respectivamente, o que, em termos A Tabela 3.12 apresenta os Índices de Facilidade e Discriminação ( ponto-bisserial ) das questões objetivas do Componente de Conhecimento Específico para os estudantes de Tecnologia em Agronegócio.
Para facilitar a diferenciação das questões, usaremos as mesmas cores da Tabela 3.10 para as diferentes classificações dos Índices de Facilidade e de Discriminação.
A partir do Índice de Facilidade obtido, pode -se concluir que quase três quartos das questões objetivas da prova foram considerados, pelo menos, Difícil : das 27 questões válidas, 20 foram classificadas como Difícil ou como Muito difícil.
Duas questões foram classificadas como Fácil, e outras cinco consideradas como Médio.
Já quanto aos índices de discriminação das questões objetivas do Componente de Conhecimento Específico da prova, tem -se como resultado a seguinte classificação : seis das 27 questões foram consideradas como boas, enquanto uma delas teve Índice de Discriminação Muito bom.
Assim, para um pouco mais de um quarto das questões – sete em 27 – os Índices de Discriminação foram Bom ou Muito bom.
Dentre as demais, 10 delas foram classificadas como Médio e outras 10 como Fraco, sendo 20, por conseguinte, a quantidade de questões nos dois patamares mais baixos de discriminação.
Constata -se, assim, que a prova – no que se refere ao Componente de Conhecimento Específico – possuía pouca capacidade de discriminar entre aqueles que dominam ou não o conteúdo.
Dentre as questões que alcançaram os maiores índices de discriminação, uma delas, a de número 32, foi classificada com Índice de Discriminação Muito bom, situando -se no intervalo de 0,40 a 0,50 do índice, e foi classificada na categoria Médio, quanto ao Índice de Facilidade.
A questão de número 30 foi a mais difícil dentre as 27 questões específicas, com baixo Índice de Facilidade, apenas 8 \% de acertos.
Essa questão apresentou poder discriminatório igualmente baixo, – 0,02, o que comprova ter sido esta questão a mais difícil para os estudantes.
Destacam -se, também, as questões 19 e 21, com Índices de Facilidade 0,11 e 012, respectivamente, o que, em termos percentuais, corresponde a 11,0 \% e 12 \% de estudantes que as responderam acertadamente.
Já 0,18 e 0,02 foram os seus Índices de Discriminação.
Tais questões foram, portanto, pelo critério ponto-bisserial, consideradas inadequadas.
Por isso, as questões 19, 21 e 30, além das questões 10, 13, 20, 23, 26, 28 e 33, foram eliminadas do cômputo da nota final.
A título de exemplo das análises do comportamento das questões objetivas, o Gráfico 3.5 analisa a questão 32 do Componente de Conhecimento Específico.
Esta questão foi considerada pelas respostas dos estudantes avaliados na prova como uma questão com nível de facilidade Médio, 0,53, ou seja, 53,0 \% dos estudantes assinalaram acertadamente a opção E, correspondente ao gabarito.
Seu Índice de Discriminação foi percentuais, corresponde a 11,0 \% e 12 \% de estudantes que as responderam acertadamente.
Já 0,18 e 0,02 foram os seus Índices de Discriminação.
Tais questões foram, portanto, pelo critério ponto-bisserial, consideradas inadequadas.
Por isso, as questões 19, 21 e 30, além das questões 10, 13, 20, 23, 26, 28 e 33, foram eliminadas do cômputo da nota final.
A título de exemplo das análises do comportamento das questões objetivas, o Gráfico 3.5 analisa a questão 32 do Componente de Conhecimento Específico.
Esta questão foi considerada pelas respostas dos estudantes avaliados na prova como uma questão com nível de facilidade Médio, 0,53, ou seja, 53,0 \% dos estudantes assinalaram acertadamente a opção E, correspondente ao gabarito.
Seu Índice de Discriminação foi igual a 0,41, classificado como Muito bom, também sendo esta questão a que apresentou o maior valor.
Neste gráfico, cada uma das cinco curvas representa o percentual de respostas em determinada alternativa da questão 32, em função do número de acertos dos estudantes nessa parte da prova, antes de possíveis eliminações de questões pelo critério do ponto-bisserial.
A alternativa correta E, representada no gráfico pela curva em vermelho, foi escolhida em maiores proporções pelos alunos com desempenho melhor nessa parte da prova.
Já as alternativas incorretas, também denominadas distratores, foram selecionadas, principalmente, por aqueles com notas mais baixas.
Neste caso também a soma não é sempre 100 \% por causa das questões não respondidas ou com mais de uma opção marcada.
Aqueles com nota zero, na sua quase totalidade deixaram esta questão em branco ou marcaram mais de uma alternativa, comportamento considerado inválido.
A proporção de alunos que selecionou a resposta correta E aumenta gradativamente, chegando a atingir 100 \% para 17 acertos, enquanto a proporção dos que escolheram alternativas incorretas decai, a partir de sete acertos, como função do número de acertos nessa parte da prova.
Os gráficos relativos às demais questões do Conhecimento Específico constam do Anexo I. 3.3 - ANÁLISE DAS QUESTÕES DISCURSIVAS Esta seção apresenta estatísticas selecionadas e histogramas das Questões Discursivas de Formação Geral ( 3.3.1 ) e Conhecimento Específico ( 3.3.2 ).
São também apresentadas e comparadas as médias de subpopulações caracterizadas por Grande Região. 3.3.1 - Componente de Formação Geral As análises dos resultados de desempenho dos estudantes de Tecnologia em Agronegócio nas duas questões discursivas relativas à Formação Geral encontram -se na Tabela 3.13 e no Gráfico 3.6.
Na Tabela 3.13, observa -se que a nota Média nesse conjunto de questões foi acima da obtida nas objetivas.
Os estudantes, de todo o Brasil, obtiveram, em Formação Geral, Média 33,5 nas questões objetivas e 51,9 nas questões discursivas.
Pode -se notar também que o Desvio padrão nesse conjunto de questões foi um pouco maior do obtido nas objetivas, 18,8 nas questões objetivas e 19,9 nas questões discursivas.
A maior Média foi obtida na região Nordeste ( 52,8 ), e a menor, na região Norte ( 48,6 ).
A Mediana de todo o Brasil, neste componente, foi 56,0, a mesma obtida nas regiões Sudeste e Centro-Oeste.
Na região Norte a Mediana foi menor, 53,0.
Nas regiões Nordeste e Sul, foi maior que o valor nacional, 56,5 e 57,5, respectivamente.
A nota Máxima do Brasil ( 94,5 ) foi obtida na região Sul, sendo 79,5 na região Norte, 86,5 a nota Máxima na região Centro-Oeste, 87,5 na região Nordeste e 88,0 na Sudeste.
A nota Mínima ( 0,0 ) foi a mesma em todas as regiões do Brasil.
O Gráfico 3.6 representa a distribuição das notas nas questões discursivas do Componente de Formação Geral.
A moda desta distribuição ocorre no intervalo ( 50 ; 60 ], seguido bem de perto pelo intervalo ( 60, 70 ].
Destaca -se também o intervalo [ 0 ; 10 ] com distribuição em torno de 7 \% do total de notas, sendo que no intervalo [ 0 ; 10 ] inclui -se, além da nota zero, a frequência de alunos que deixaram este tipo de questão em branco.
Na sequência, os resultados verificados para cada uma das questões discursivas de Formação Geral serão apresentados, estabelecendo -se relações com os conteúdos abordados em cada uma delas.
Os comentários da Banca de docentes corretores, a respeito do observado na correção das respostas dos estudantes, suas impressões e conclusões serão apresentados junto à análise de cada questão.
Cumpre esclarecer que, tendo em vista que as questões discursivas de Formação Geral são padronizadas, ou seja, constam de todas as provas, os comentários da Banca são os mesmos para todas as carreiras acadêmicas, sendo direcionados a todos os estudantes que participaram do Enade/2016.
A seguir, serão analisados os desempenhos da Área de Tecnologia em Agronegócio nas duas questões discursivas de Formação Geral do Enade/2016, comparando -se os resultados obtidos com comentários para cada questão. 3.3.1.1 - Análise de Conteúdo da Questão Discursiva 1 do Componente de Formação Geral Os dados de Tecnologia em Agronegócio, obtidos a partir das respostas à questão 1, encontram -se na Tabela 3.14 e no Gráfico 3.7.
Nessa questão – de melhor desempenho dentre as duas de Formação Geral – os alunos, de todo o Brasil, tiveram Média 58,1.
A maior Média para a questão 1 foi obtida na região Centro-Oeste ( 58,8 ), e a menor, na região Nordeste ( 54,9 ).
Quanto à variabilidade das notas, o Desvio padrão de todo o Brasil foi 28,9.
O menor Desvio padrão foi obtido na região Sudeste ( 28,1 ), e o maior Desvio padrão foi obtido na região Nordeste ( 30,2 ).
As medianas de todas as regiões e a do Brasil como um todo, foram iguais ( 60,0 ).
As notas máximas e mínimas da questão discursiva 1 foram, sem exceção, as mesmas para todas as regiões do Brasil, respectivamente, 100,0 e 0,0.
O Gráfico 3.7 mostra a distribuição das notas na questão discursiva 1 do Componente de Formação Geral.
Observa -se que a maior frequência corresponde aos alunos que obtiveram nota dentro do intervalo ( 70 ; 80 ], correspondendo à moda da distribuição.
Destacam -se, também, os alunos que deixaram a questão em branco com aproximadamente 9 \% do total. 3.3.1.2 - Comentários sobre a correção de Conteúdo das respostas à Questão Discursiva 1 A questão discursiva 1 alinha -se aos conteúdos programáticos previstos nas diretrizes curriculares, quanto às habilidades e às competências para a atualização permanente e aos conhecimentos sobre a realidade brasileira, mundial, tangenciando outras áreas do conhecimento.
O enunciado é claro e indica um tema de extrema relevância, qual seja, as migrações humanas e todo o complexo de problemas que esse movimento causa, quer na origem, quer no seu ponto de destino.
A questão é extremamente atual, maximizada pelos acompanhamentos midiáticos que espetacularizam cenas dramáticas do processo de fuga de massas humanas das guerras, como a morte de imigrantes.
Além do aspecto internacional, deve se ter em conta que o Brasil sofreu vários influxos migratórios, desde o forçado, pela transferência de escravos da África, até o incentivado, dos italianos, alemães e outros europeus para trabalhar na agricultura.
Isso sem considerar a própria \textbf{colonização} portuguesa originária e as migrações posteriores dos mesmos portugueses.
Posteriormente, o Brasil tornou -se um país de emigrantes que foram aportar a outros locais na busca de melhor situação econômica.
Muitos migraram para a América do Norte e para a Europa, principalmente.
Mais recentemente, por força do crescimento econômico pátrio, novas levas de imigrantes desembarcaram no Brasil, muitos deles provenientes do Haiti, embora tenhamos recebidos vizinhos latinos, como bolivianos, paraguaios, chilenos, argentinos, venezuelanos, peruanos, colombianos que também vieram em busca de uma vida melhor.
Como se verifica, a temática indicada no enunciado pertence à realidade conhecida amplamente pelos estudantes, pois, também divulgada por todos os meios de comunicação.
Trata -se de tema geral que, normalmente, não se inclui nas grades curriculares superiores, mas é de conhecimento geral e relacionado aos direitos humanos, este sim, ponto de integração às diretrizes curriculares.
A linguagem utilizada foi clara, indicando a existência dos fatos vinculados aos movimentos migratórios.
Por outro lado, o texto demanda uma reflexão sobre a realidade dos indivíduos que são compelidos ao afastamento dos seus lares por motivos externos, como guerras, conflitos civis, religiosos, perseguições das mais variadas, questões econômicas.
Como se trata de texto contendo uma ampla gama de possibilidades, é natural que ocorram interpretações diversas das esperadas.
Do mesmo modo, a contemporaneidade do tema encaminha o enunciado para ser qualificado como questão de média dificuldade, em comparação com as demais questões apresentadas.
O padrão de resposta indicou possíveis interpretações do texto.
No que concerne ao grau de profundidade, mostrou a necessidade de conhecimento de vários itens pertinentes à migração e que são refletidos nas várias possibilidades de resposta previstas pelo padrão.
Pela análise das provas realizadas, constatou -se que muitos foram os que identificaram os itens indicados no padrão de resposta.
Essa circunstância poderia ser justificada, tendo em vista que o tema vem sendo bastante difundido na mídia social, gerando elementos naturais para a resposta, bem como da história do país, que, intimamente, vincula -se ao processo migratório, quer ativo, remetendo pessoas para outros países, quer passivo, na recepção de pessoas.
Essa característica também permite responder aos termos postos quanto aos Estados que recebem os migrantes.
Assim, a indicação dos aspectos positivos e negativos também é bastante difundida.
As respostas medianas e boas apreenderam o sentido do enunciado, e indicaram os motivos expressos no padrão de respostas, como os ensejadores das imigrações, notadamente, para a Europa, que seria o foco do texto.
Por outro lado, também identificaram os problemas dos imigrantes, no mesmo padrão indicado para as respostas desejadas.
Alguns concluintes procuraram exemplos de comparação na realidade brasileira, como os clássicos imigrantes italianos, alemães e de outras nacionalidades, bem como os atuais haitianos e demais povos da América Latina.
Uma minoria não apreendeu o sentido do enunciado e, somente, apresentou elementos da realidade nacional.
Esses foram avaliados no conjunto de provas fracas.
As respostas demonstraram que o concluinte está atualizado, em regra, com o tema da imigração, indicando os fundamentos dessa movimentação populacional que, nos últimos tempos, tem se agravado pela incidência de conflitos armados em vários países, ensejadores de reflexos na segurança dos habitantes desses locais, bem como o surgimento de focos de doenças e ausência dos serviços públicos básicos, além da fome, decorrente da impossibilidade de produzir alimentos ou mesmo comerciá -los.
Do ponto de vista dos Estados que recebem o fluxo de imigrantes, também restou clara a percepção das dificuldades vinculadas à ausência de estrutura quanto aos serviços públicos ( saúde, educação, segurança, moradia, alimentação ) para receber um número expressivo de imigrantes.
Em época de recessão econômica, a presença de estrangeiros acarreta a competição pela busca de empregos, gera movimentos xenofóbicos, bem como choques culturais e políticos.
Deve -se ter em mira que esses movimentos têm um viés de insegurança, pelos países receptores, diante da crise causada por movimentos considerados terroristas, embora a imensa maioria não tenha qualquer vínculo com esses movimentos.
Os recentes atentados na Europa indicam que os seus autores têm origem em comunidades de imigrantes, embora muitos já tenham nascido nos países da comunidade europeia.
Essas características negativas, no entanto, paradoxalmente, convivem com aspectos positivos, o que encaminha a questão para a necessidade dos jovens 3.3.1.3 - Análise de Conteúdo da Questão Discursiva 2 do Componente de Formação Geral A Tabela 3.15 mostra que o desempenho médio dos estudantes na questão discursiva 2 ( média 43,4 ) foi inferior ao obtido na questão discursiva 1 ( média 58,1 ).
A região Nordeste foi aquela cuja Média foi maior ( 46,9 ), e a de menor Média foi a região Norte ( 39,2 ).
Quanto à variabilidade das notas, o Desvio padrão de todo o Brasil foi 22,8, inferior ao obtido na questão discursiva 1 ( 28,9 ).
O maior desvio nessa questão foi obtido na região Norte ( 24,4 ), enquanto o menor, foi obtido na região Nordeste ( 22,3 ).
A Mediana de todo o Brasil foi 50,0, a mesma de todas as regiões.
A nota Máxima ( 100,0 ) foi obtida nas regiões Nordeste, Sudeste e Sul.
Na região Norte, a nota Máxima foi 80,0, e na região Centro-Oeste foi 90,0.
A nota Mínima foi 0,0 e ocorreu em todas as regiões, sem exceção.
O Gráfico 3.8 mostra a distribuição das notas na questão discursiva 2 do Componente de Formação Geral.
Observa -se que a maior frequência corresponde aos alunos que obtiveram nota dentro do intervalo ( 50 ; 60 ], com quase 25 \% do total das respostas.
Destaca -se, também, o grupo dos alunos que deixaram a questão em branco, caracterizando -se como máximo local.
Nota -se, ainda, que as notas ficaram mais imigrantes, que ocuparão postos de trabalho não desejados pelos habitantes locais, bem como, compensarão os baixos índices de crescimento populacional.
Existe outro reflexo positivo que está relacionado ao sistema previdenciário, que terá novos contribuintes nos trabalhadores estrangeiros que se integram ao sistema.
Esses pontos foram abordados e indicam que existe uma consciência dos concluintes dos cursos universitários quanto aos principais itens relacionados ao tema que foi posto para análise. 3.3.1.3 - Análise de Conteúdo da Questão Discursiva 2 do Componente de Formação Geral A Tabela 3.15 mostra que o desempenho médio dos estudantes na questão discursiva 2 ( média 43,4 ) foi inferior ao obtido na questão discursiva 1 ( média 58,1 ).
A região Nordeste foi aquela cuja Média foi maior ( 46,9 ), e a de menor Média foi a região Norte ( 39,2 ).
Quanto à variabilidade das notas, o Desvio padrão de todo o Brasil foi 22,8, inferior ao obtido na questão discursiva 1 ( 28,9 ).
O maior desvio nessa questão foi obtido na região Norte ( 24,4 ), enquanto o menor, foi obtido na região Nordeste ( 22,3 ).
A Mediana de todo o Brasil foi 50,0, a mesma de todas as regiões.
A nota Máxima ( 100,0 ) foi obtida nas regiões Nordeste, Sudeste e Sul.
Na região Norte, a nota Máxima foi 80,0, e na região Centro-Oeste foi 90,0.
A nota Mínima foi 0,0 e ocorreu em todas as regiões, sem exceção.
O Gráfico 3.8 mostra a distribuição das notas na questão discursiva 2 do Componente de Formação Geral.
Observa -se que a maior frequência corresponde aos alunos que obtiveram nota dentro do intervalo ( 50 ; 60 ], com quase 25 \% do total das respostas.
Destaca -se, também, o grupo dos alunos que deixaram a questão em branco, caracterizando -se como máximo local.
Nota -se, ainda, que as notas ficaram mais 3.3.1.4 - Comentários sobre a correção de Conteúdo das respostas à Questão Discursiva 2 A questão discursiva 2 tem seu conteúdo integrado às Diretrizes Curriculares, uma vez que a proteção aos direitos humanos e às relações de gênero constituem matrizes fundamentais para conhecimento dos profissionais de nível superior.
O enunciado é claro e indica um tema de extrema relevância, qual seja, a análise da questão da mulher submetida a injustas agressões e os seus reflexos pessoais, sociais e familiares.
Também enfrenta a busca pela maior proteção no concernente às agressões, com a instituição de um modelo legal mais repressivo e que foi gerado por uma situação especifica de ataque à mulher, com consequências dramáticas para a vítima, Maria da Penha, cuja lei levou seu nome.
Deve ser recordado que os conflitos vinculados às relações familiares sempre foram circunscritos ao seu âmbito e, quando dele extravasavam, por força de agressões contundentes à integridade física de um dos partícipes da relação, em geral a mulher, raramente redundavam em punições relevantes, sendo comum o pagamento em cestas básicas de alimentos, como sanção.
Após a edição da lei Maria da Penha ( Lei no. 11.340, sancionada em 7 de agosto de 2006 ), além das sanções sofrerem modificações, houve a criação de modelos de proteção e atendimento especial para os casos de violência doméstica, procurando proteger as mulheres agredidas antes de danos maiores, muitas vezes irreversíveis, incluindo aí o feminicídio, também incluído no Código Penal, como circunstância qualificadora do crime de homicídio, e incluí -lo no rol dos crimes hediondos.
A lei Maria da Penha busca realizar, como indicado no seu pórtico, a Convenção sobre a Eliminação de Todas as Formas de Discriminação contra as Mulheres e a Convenção Interamericana para Prevenir, Punir e Erradicar a Violência contra a Mulher.
No seu texto, várias são as formas de violência doméstica e familiar contra a mulher, identificadas como sendo, entre outras : “ a violência física, entendida como qualquer conduta que ofenda sua integridade ou saúde corporal ; a violência psicológica, entendida como qualquer conduta que lhe cause dano emocional e diminuição da autoestima ou que lhe prejudique e perturbe o pleno desenvolvimento ou que vise a degradar ou controlar suas ações, comportamentos, crenças e decisões, mediante ameaça, constrangimento, humilhação, manipulação, isolamento, vigilância constante, perseguição contumaz, insulto, chantagem, ridicularização, exploração e limitação do direito de ir e vir ou qualquer outro meio que lhe cause prejuízo à saúde psicológica e à autodeterminação ; a violência sexual, entendida como qualquer conduta que a constranja a presenciar, a manter ou a participar de relação sexual não desejada, mediante intimidação, ameaça, coação ou uso da força ; que a induza a comercializar ou a utilizar, de qualquer modo, a sua sexualidade, que a impeça de usar qualquer método contraceptivo ou que a force ao matrimônio, à gravidez, ao aborto ou à prostituição, mediante coação, chantagem, suborno ou manipulação ; ou que limite ou anule o exercício de seus direitos sexuais e reprodutivos ; a violência patrimonial, entendida como qualquer conduta que configure retenção, subtração, destruição parcial ou total de seus objetos, instrumentos de trabalho, documentos pessoais, bens, valores e direitos ou recursos econômicos, incluindo os destinados a satisfazer suas necessidades ; a violência moral, entendida como qualquer conduta que configure calúnia, difamação ou injúria. ” ( Conforme texto extraído do artigo sétimo da referida lei ).
Também se buscam medidas de integração entre os vários órgãos públicos, abrangendo repressão, proteção, conscientização, bem como utilização de equipes multidisciplinares, tendo em vista os graves danos, de variados matizes, que redundam das práticas agressivas.
Nessa linha, também existe na lei a previsão de locais de atendimento específico ( centros de atendimento integral e multidisciplinar para mulheres e respectivos dependentes em situação de violência doméstica e familiar ; casas-abrigo para mulheres e respectivos dependentes menores em situação de violência doméstica e familiar ; delegacias, núcleos de defensoria pública, serviços de saúde e centros de perícia médico-legal especializados no atendimento à mulher em situação de violência doméstica e familiar ; programas e campanhas de enfrentamento da violência doméstica e familiar ; centros de educação e de reabilitação para os agressores ).
O Poder Judiciário, também, deve especializar o seu atendimento para aplicar as medidas de restrição ao agressor de forma mais eficiente.
A linguagem utilizada na questão foi clara, indicando a grave violação dos direitos das mulheres injustamente agredidas e a luta pela busca de sanções adequadas aos agressores e de proteção pessoal e para os de sua família.
Ademais, o tema não está circunscrito ao Brasil, mas o país possui um número muito expressivo de atos que compõem esse cenário.
Por outro lado, o texto demanda uma reflexão sobre a realidade, e abarca diversos segmentos da sociedade, pois as situações de agressão não se restringem às camadas mais populares.
Como se trata de texto contendo uma ampla gama de possibilidades, é natural que ocorram interpretações diversas das esperadas.
O padrão de resposta identifica a projeção do elaborador na resposta a ser apresentada pelo estudante.
Está em linguagem correta e indica possíveis interpretações do texto.
Suas opções foram bastante amplas, abrangendo a maior parte das respostas encontradas.
No que concerne ao grau de profundidade, o padrão de resposta indicou a necessidade de conhecimento de lei especial que regula a violência doméstica, o que exige ou qualificação da área de Direito ou que os currículos dos cursos dos participantes desse ciclo do Enade tenham oferecido alguma disciplina introdutória da área jurídica.
Ao analisar as provas realizadas, constatamos que, na primeira parte, houve maior facilidade em responder aos itens indicados no padrão de resposta, pois atinentes, não só à formação profissional, mas, também, ao conhecimento da realidade, pois, infelizmente, a violência doméstica está disseminada na sociedade e nos vários pontos do país.
Mas, verificou -se que as respostas vieram, no mais das vezes, desorganizadas, dificultando a correção.
Apesar disso, restou claro que a maioria esmagadora das respostas confluiu para as indicações do padrão apresentado, sendo relevante indicar que o tema remeteu para a luta em relação à igualdade de gênero, temática candente em especial para as concluintes do sexo feminino.
Diversos textos se apresentaram de forma concatenada, fundamentada e, muitas vezes, lamentando que, nos dias atuais, ainda se discuta a necessidade de proteger as mulheres dessa espécie de violência.
Em muitos casos, também, os concluintes fizeram um liame entre o acesso à educação e a melhora das relações entre homens e mulheres, ao permitir que a conclusão de maior grau de educação formal conduzisse à diminuição do índice de violência contra a mulher.
Foram identificadas respostas consoantes com o padrão apresentado, indicando os itens que nele constam em relação aos reflexos pessoais, na família e na sociedade.
As respostas dos estudantes raramente faziam referência à organização protetiva instituída pela Lei Maria da Penha, como esperado pelo padrão de resposta.
A respeito de tal aspecto o padrão indicava como uma das possibilidades de respostas ( item ‘ b ’ ) : a criação de casas de refúgio ou casa-abrigo, para acolher mulheres vítimas de violência doméstica ; reforço às Delegacias de Atendimento à Mulher ; inclusão da vítima em programas assistenciais do governo, programas de proteção à vítima e à testemunha, transferência de local de trabalho, se a vítima for servidora pública.
Como paradigma, indicaram -se as Delegacias da Mulher que são organizadas, com a quase integralidade dos seus membros de policiais do sexo feminino, e especializadas no atendimento a mulheres submetidas a situações de violência.
Essa preocupação gerencial tem por escopo bloquear os preconceitos que acompanham essas dramáticas histórias de violência, o que gera uma incompreensão e uma ineficiência na rede de proteção.
As respostas demonstram que os concluintes têm ciência do problema, da sua gravidade, do vínculo com a questão de gênero e com a submissão da mulher como efeito da sociedade patriarcal, ainda que tratado de forma superficial, com base na mídia e, algumas vezes, com apoio de experiências pessoais.
Identificou -se uma generalizada compreensão de que o problema está sendo enfrentado e que ainda existe um longo caminho no percurso para a sua resolução. 3.3.1.5 - Análise de Língua Portuguesa das Questões Discursivas do Componente de Formação Geral Os dados de Tecnologia em Agronegócio, obtidos a partir das respostas às questões discursivas do Componente de Formação Geral, no que tange à Língua Portuguesa, encontram -se na Tabela 3.16 e no Gráfico 3.9.
Nesse aspecto, os alunos, de todo o Brasil, tiveram Média 56,8.
A maior Média com respeito à Língua Portuguesa foi obtida na região Sudeste ( 57,9 ), e a menor, na região Norte ( 53,3 ).
Quanto à variabilidade das notas, o Desvio padrão de todo o Brasil foi 18,6.
O menor Desvio padrão foi obtido na região Nordeste ( 16,9 ) e o maior Desvio padrão foi obtido na região Norte ( 20,3 ).
A Mediana das notas de Língua Portuguesa foi 62,5 para o Brasil como um todo.
A maior Mediana com respeito à Língua Portuguesa foi obtida nas regiões Sudeste e Sul ( 62,5 ), e a menor, na região Norte ( 57,5 ).
A nota Máxima para todo o Brasil foi de 87,5, com, pelo menos, um aluno tirando essa nota na região Sul.
Além disso, a nota Mínima foi zero em todas as regiões do país.
O Gráfico 3.9 mostra a distribuição das notas de Língua Portuguesa do Componente de Formação Geral.
Observa -se que a maior frequência ( quase 35 \% ) corresponde aos alunos que obtiveram nota no intervalo ( 60 ; 70 ].
Destacam -se, também, os alunos que deixaram ambas as questões em branco, representando mais de 5 \% do total, caracterizando -se como máximo local. 3.3.1.6 - Comentários sobre a correção das respostas de Formação Geral com respeito à Língua Portuguesa Na avaliação do desempenho linguístico dos participantes foram consideradas quatro competências distintas, de modo a permitir um mapeamento detalhado do domínio dos recursos disponíveis na Língua Portuguesa para a comunicação escrita formal.
Nessa situação comunicativa, as exigências para o desenvolvimento desse tipo de texto dizem respeito : à adequação da seleção vocabular, ao desenvolvimento do conteúdo, à estruturação sintática dos períodos, à organização lógica das ideias, à utilização de procedimentos de encadeamento textual e de referenciação, à obediência às exigências morfossintáticas próprias da modalidade escrita da norma-padrão, ao respeito às regras ortográficas e às regras de acentuação gráfica.
A avaliação foi realizada considerando -se os textos das respostas às duas questões de Formação Geral e os resultados obtidos em cada uma delas, como seria de se esperar, foram semelhantes.
As competências avaliadas podem ser sintetizadas nos aspectos descritos a seguir.
Aspectos textuais - estruturação textual condizente com o gênero discursivo e o modo de organização textual expositivo/argumentativo adequado ao gênero – essa competência envolve : a estruturação sintática condizente com o padrão da modalidade escrita formal da Língua Portuguesa de modo a garantir a clareza necessária ; a distribuição do conteúdo do texto em parágrafos, de modo a garantir a sua organização temática ; a utilização de operadores discursivos que contribuam para a progressão temática do texto, estabelecendo relações lógicas entre as ideias apresentadas, tanto do ponto de vista intrafrasal, como do interfrasal ; a utilização de procedimentos de referenciação lexical e pronominal que permitam a retomada de referentes textuais ; o respeito às regras de pontuação como fator de estruturação do período.
Aspectos morfossintáticos - domínio dos diferentes aspectos morfossintáticos próprios da modalidade escrita formal da norma-padrão da Língua Portuguesa – essa competência envolve : a concordância nominal, a concordância verbal, a regência nominal, a regência verbal, a flexão nominal, a flexão verbal, a correlação entre os tempos verbais, a colocação pronominal e a utilização de sinais de pontuação que contribuam para a organização lógica da frase e do texto.
Aspectos vocabulares - seleção vocabular adequada à modalidade escrita formal da Língua Portuguesa, exigida pela situação comunicativa – essa competência envolve a precisão na seleção/utilização do vocabulário relacionado à temática solicitada pela questão ; a ausência de marcas da oralidade, como termos de sentido muito genérico ( “ coisa ”, “ negócio ”, “ você ” ) e termos de registros mais informais ( como gírias, jargões, frases feitas, ditados populares, termos regionais ).
Assim, espera -se que o participante respeite a adequação vocabular não usando gírias ou expressões coloquiais, evite repetição desnecessária de palavras e utilize um vocabulário mais formal, como solicitado por um texto dissertativo.
Aspectos ortográficos - respeito às convenções ortográficas da norma-padrão da Língua Portuguesa – essa competência envolve o domínio das regras de acentuação gráfica e da grafia padrão das palavras ( com ausência de abreviaturas próprias da linguagem da internet ), de acordo com as convenções estabelecidas pela legislação em vigor e consubstanciadas no Vocabulário Ortográfico da Língua Portuguesa, editado pela Academia Brasileira de Letras. ( Vale notar que, nessa edição de 2016, não será aceita a legislação anterior, no caso das regras relativas ao uso do hífen e da acentuação gráfica ).
A escolha dessas competências para subsidiar o processo de avaliação apoia- apresentado uma intensa simplificação, originada no padrão da modalidade oral da Língua Portuguesa.
No caso do texto de base dissertativa, inscrito em um registro formal, a distância entre as duas modalidades é ainda maior, o que provoca situações de hipercorreção ( desvios provocados pela incorporação indevida de uma regra da norma-padrão ) e de truncamentos sintáticos ( estruturas frasais incompreensíveis devido à complexidade sintática própria da modalidade escrita ).
Para efeito de pontuação, o critério aprovado pelo Inep, e já utilizado na avaliação do desempenho linguístico nas edições do Enade desde 2013, reúne as competências acima descritas em três grupos da seguinte forma : aspectos textuais ( 40 \% ), aspectos morfossintáticos e vocabulares ( 40 \% ) e aspectos ortográficos ( 20 \% ).
A seguir, analisam -se qualitativamente os resultados da avaliação de cada um dos aspectos avaliados nos textos de resposta das questões discursivas 1 e 2 de Formação Geral.
Aspectos textuais Esta competência é a que se revelou como a mais problemática entre os participantes, porque são muitos os desvios observados.
Em relação à estrutura formal do texto produzido pelos participantes, observaram -se relações linguísticas quase agramaticais, como as estabelecidas pela sequência de gerúndios sem o apoio de um ponto de partida para a organização das informações gramaticais e semânticas, ou seja, sem oração principal.
Observou -se que uma grande parte dos participantes não organizou as ideias em parágrafos, talvez devido ao pequeno número de linhas disponibilizadas para a resposta da questão ou, quem sabe, pela suposição de que não seria necessária essa divisão por não se tratar de um texto no modelo de uma redação dissertativo-argumentativa, como solicitado nos vestibulares.
Quanto à utilização dos mecanismos de referenciação, deve -se destacar a ocorrência, em uma parte dos textos, de repetições de palavras ou expressões sem a utilização de termos sinônimos ou pronomes, como seria adequado.
Outro aspecto relevante a ser destacado é a quase total ausência de operadores argumentativos, tanto intrafrasais como interfrasais, repercutindo uma tendência atual da mídia escrita.
A respeito do uso dos sinais de pontuação, observaram -se vários desvios nos textos analisados, com a ocorrência, inclusive, de parágrafos sem marca interna de pontuação para separar os períodos.
São os seguintes os tipos de problemas encontrados, que se caracterizam como desvios ao estabelecido pela norma-padrão : a ) uso inadequado de vírgula : • ocorrência entre o sujeito e o predicado, o que é considerado inadequado por interromper a sequência natural da frase ; • ocorrência de apenas uma das vírgulas para intercalar uma palavra, uma expressão ou uma oração encaixada, quando o correto seria a demarcação com uma vírgula antes e uma depois do termo encaixado na frase ; • ocorrência da vírgula no lugar do ponto para separar ideias que constituem períodos distintos ; b ) uso inadequado de ponto e vírgula para separar elementos que deveriam ser separados por vírgula, por estarem dentro do mesmo período ; c ) ausência de ponto final para separar períodos.
Vale observar que não foi penalizada a ausência de vírgula para destacar locuções ou adjuntos adverbiais de pequena extensão deslocados de posição na frase, por ser um uso opcional.
Os problemas relativos à estruturação textual geram, portanto, os seguintes tipos de desvios : • sequência justaposta de ideias sem encaixamentos sintáticos ; • redução de estruturas subordinadas, ao lado do aumento na frequência de estruturas coordenadas e absolutas ; • redução no uso de conectores para expressar relações lógicas essenciais à construção do texto ; • emprego inadequado de operadores que não estabelecem relações lógicas coerentes entre ideias do texto ; • repetição de termos sem a utilização de procedimentos mais sofisticados de substituição ( hiperonímias, hiponímias, nominalizações, expressões metafóricas ) ; • frases fragmentadas que comprometem a estrutura lógico-gramatical ; • frases formadas apenas por oração subordinada, sem oração principal.
Aspectos morfossintáticos e vocabulares Os resultados relativos aos aspectos morfossintáticos ( concordância, regência, colocação de pronomes, flexão nominal e verbal ) evidenciam desvios recorrentes no desempenho dos participantes, relacionados ao seu processo de formação e, principalmente, a hábitos da oralidade, marcados pela informalidade.
Apesar de exemplificarem processos de mudança em curso na Língua Portuguesa, esses desvios foram penalizados por não atenderem às exigências da norma-padrão relativas ao texto escrito em situação formal, como deveria ser o texto dissertativo solicitado no Enade.
O desvio mais frequente, em relação à regência, é a falta do sinal indicativo da crase – isso revela que o usuário não tem consciência de que, sob a forma do termo “ a ”, existe a presença de uma preposição “ a ”, existe a presença de uma preposição “ a ”, exigida pela regência do termo anterior, combinada a um artigo definido ( por exemplo : “ Devido as grandes guerras que estão acontecendo... ” ) Outros problemas relacionados à regência verbal e à nominal foram identificados em alguns textos dos participantes : ( a ) a ausência de preposição ou seu uso indevido antes de pronome relativo ( por exemplo : “ a região que eles moram não oferece condições de segurança ” ) ; ( b ) o uso indevido do pronome relativo ( por exemplo : “ não pensam na cidade para que eles vão ” ) ; ( c ) utilização de pronome “ onde ” deslocado de sua função locativa ).
Esses processos refletem hábitos da modalidade oral da língua, em situações de registro informal.
Apesar da possibilidade de que essa alteração de regência se generalize no padrão escrito da Língua Portuguesa, como já está ocorrendo até em textos jornalísticos, a ausência da preposição foi penalizada neste processo de avaliação.
Outro desvio observado nos textos diz respeito aos processos de concordância verbal e de concordância nominal.
Quanto à concordância de número, além das ocorrências já conhecidas de falta da marca de plural em verbos ou adjetivos, em função da ocorrência de substantivo plural na função de sujeito, observaram -se dois tipos de desvio já registrados nas pesquisas linguísticas : ( a ) a ausência de marca de plural com sujeito posposto ( por exemplo : “ Foi expulso do país os imigrantes que entraram sem autorização da polícia ” ) ; ( b ) o uso indevido da marca de plural para estabelecer concordância com o plural da locução adjetiva, apesar de o substantivo que funciona como núcleo do sintagma nominal estar no singular ( por exemplo : “ A proteção das mulheres que sofrem violência foram garantidas pela Lei Maria da Penha ” ).
Uma ocorrência generalizada foi a ausência de acento circunflexo na forma plural do presente do indicativo do verbo “ ter ”, que foi considerada como um desvio na concordância verbal e não meramente de acentuação gráfica.
Quanto à concordância de gênero, alguns casos foram observados, normalmente no âmbito de sintagmas nominais longos, em que o adjetivo está afastado do substantivo.
Quanto à questão da colocação pronominal, foram poucos os desvios observados.
Apesar de serem aspectos relacionados à oralidade, observou -se que a maioria dos participantes já incorporou regras como a não introdução da frase por um pronome oblíquo e a próclise na presença de um termo atrator.
Não se adotou, entretanto, o padrão excessivamente formal descrito pelas gramáticas normativas em relação à posição do pronome oblíquo em locuções verbais, já que esse uso está muito distante da prática cotidiana, até em textos mais formais. relação à posição do pronome oblíquo em locuções verbais, já que esse uso está muito distante da prática cotidiana, até em textos mais formais.
Quanto aos aspectos vocabulares ( precisão semântica do vocabulário utilizado e adequação ao registro formal exigido pelo tipo de texto solicitado no enunciado ), podem ser destacados os seguintes tipos de desvios : expressões da oralidade, seleção vocabular incompatível com o contexto, gerando situações de falta de inteligibilidade ; falta de domínio de vocabulário mais abstrato e de maior complexidade, essencial ao desenvolvimento do texto de base dissertativa.
Não se observou presença de palavras de baixa calão, ao contrário do exame de 2015, em que o tema “ funk ” gerou a utilização desse tipo de vocabulário.
Várias marcas de oralidade foram identificadas, embora não em alta frequência : o uso do pronome relativo “ onde ” como relativo universal, falta de artigo definido antes de substantivo, repetição de palavras por falta de vocabulário, reduções como “ tá ”, “ pra ”, “ pro ”, “ prum ”, expressões informais.
Aspectos ortográficos : O desempenho dos participantes revelou uma diferença muito grande nos dois aspectos de caráter ortográfico analisados neste item : grafia e acentuação.
Observou -se a ocorrência de baixo índice de desvios na grafia das palavras e grande índice de desvios no uso dos sinais de acentuação gráfica.
Em vários casos, ocorre ausência completa de acentuação.
Os resultados revelam que existe uma tendência crescente entre os universitários brasileiros de eliminação da acentuação gráfica, talvez motivada pelos hábitos relacionados às redes sociais ou devido a um conhecimento ainda superficial sobre o Acordo Ortográfico da Língua Portuguesa de 1990.
Os casos mais sistemáticos de eliminação do acento indicador da sílaba tônica são : • palavras proparoxítonas ( por exemplo : “ proximos ”, “ politicos ”, “ publica ”, “ fisica ”, “ psicologica ” ) ; • palavras paroxítonas terminadas em ditongo crescente ( por exemplo : “ necessario ”, “ noticiarios ”, “ individuo ”, “ dependencia ”, “ varios ”, “ propria ”, “ violencia ”, “ ocorrencia ”, “ suicidio ” ) ; • palavras oxítonas ( por exemplo : “ ate ”, “ tambem ”, “ mes ” ) ; • palavras paroxítonas terminadas em ditongo nasal ( por exemplo : “ orgão ”, “ orfão ” ).
Quanto ao domínio das convenções relativas à grafia das palavras, pode -se concluir que os desvios estão majoritariamente relacionados a questões relativas à oralidade : • hipercorreção pela escolha de “ e ” no lugar de “ i ” ou de “ i ” no lugar de “ e ”, por influência de hábitos da oralidade : por exemplo, “ inviado ” ( no lugar de “ enviado ” ), “ entervenção ” ( no lugar de “ intervenção ” ), “ empedimento ” ( no lugar de “ impedimento ) ; • eliminação do “ r ” marcador do infinitivo verbal : por exemplo, “ esta ” ( no lugar de “ estar ” ), “ estuda ” ( no lugar de “ estudar ” ), “ migra ” ( no lugar de “ migrar ” ) ; • utilização de “ s ” no lugar de “ c ” : por exemplo, “ mensionado ” ( no lugar de “ mencionado ” ) ; “ insentivou ” ( no lugar de “ incentivou ” ) ; • utilização de “ s ” no lugar de “ x ” : por exemplo, “ espectativa ” ( no lugar de “ expectativa ”, devido à confusão com o verbo “ esperar ” ) ; e a situação contrária, “ x ” no lugar de “ s ”, “ expantoso ” ( no lugar de “ espantoso ” ) ; • utilização de “ ss ” no lugar de “ ç ” : por exemplo, “ gerassões ” ( no lugar de “ gerações ” ) ; • uso de “ am ” no lugar de “ ão ” na forma de terceira pessoa do singular do presente e do pretérito perfeito do indicativo : por exemplo, “ estam ” ( no lugar de “ estão ” ) ; • utilização de “ l ” no lugar de “ u ” em final de sílaba : por exemplo, “ tralma ” ( no lugar de “ trauma ” ) ; • outros casos de desvios de grafia relacionados à variação diastrática podem ser observados em “ estrupo ” ( no lugar de “ estupro ” ), “ subjulgadas ” ( no lugar de “ subjugadas ” ), “ precoseito ” ( no lugar de “ preconceito ” ), “ soubre ” ( no lugar de “ soube ” ), “ vecendo ” ( no lugar de “ vencendo ” ), “ indesencia ” ( no lugar de “ indecência ” ), “ apolojia ” ( no lugar de “ apologia ” ), “ fulga ” ( no lugar de “ fuga ” ). • inadequação no uso da maiúscula para destacar determinadas palavras-chave do texto, como “ Violência ”, “ Brasileiros ”.
Destaque -se, também, a ocorrência de participantes que grafam os textos inteiramente em caixa alta.
Vale observar que, ao contrário do esperado, não apareceram abreviaturas próprias do “ internetês ”, relacionadas ao uso de redes sociais e de e-mails.
Sua influência, portanto, parece se restringir à ausência dos acentos gráficos.
Em relação aos aspectos textuais, observa -se dificuldade na estruturação textual, o que resulta em truncamentos sintáticos e prejuízo da compreensão do texto, e desrespeito a regras essenciais de pontuação, principalmente pela dificuldade em identificar as relações lógicas que se estabelecem entre as unidades frasais.
Em relação aos aspectos morfossintáticos, domínio das regras de concordância nominal e verbal, mas desrespeito às regras de regência nominal e verbal, principalmente no que diz respeito ao uso do sinal indicativo da crase.
Em relação aos aspectos vocabulares, reduzida utilização de palavras de uso mais formal e abstrato, adequado à construção do texto de base dissertativa.
Em relação às convenções ortográficas, domínio da grafia das palavras, mas falta de atendimento à maioria das regras de acentuação gráfica. 3.3.2 - Componente de Conhecimento Específico Na parte da prova relativa às questões discursivas do Componente de Conhecimento Específico ( Tabela 3.17 ), observa -se que a Média foi bem mais baixa do que para as questões discursivas do Componente de Formação Geral.
Enquanto no Componente de Formação Geral a Média para estudantes de Tecnologia em Agronegócio de todo o Brasil foi 51,9, na parte de Conhecimento Específico a Média foi 44,1.
A maior Média deste componente foi obtida pelos estudantes da região Sul ( 48,3 ), e a menor, pelos da região Norte ( 40,7 ).
Quanto à variabilidade das notas, o Desvio padrão de todo o Brasil foi 18,4.
O maior Desvio padrão foi encontrado na região Sul ( 18,7 ), e o menor, na região Nordeste ( 17,2 ).
A maior nota Máxima, 100, foi obtida na região Sul.
Nas demais regiões a nota Máxima foi : 80,0 na região Nordeste, 83,3 na Norte, 90,0 na região Centro-Oeste e 93,3 na região Sudeste.
Além disso, a nota Mínima ( 0,0 ) foi obtida por pelo menos um aluno em todas as regiões do Brasil.
A Mediana do Brasil como um todo foi 45,0.
A maior Mediana foi obtida na região Sul ( 50,0 ), e a menor, na região Nordeste ( 40,0 ).
O Gráfico 3.10 representa a distribuição das notas nas questões discursivas no Componente de Conhecimento Específico.
Esta é uma distribuição unimodal com a moda ocorrendo no intervalo ( 40, 50 ].
A análise de cada uma destas questões será feita a seguir. 3.3.2.1 - Análise da Questão Discursiva 3 do Componente de Conhecimento Específico Na questão 3, cujos resultados aferidos encontram -se descritos na Tabela 3.18, a Média dos estudantes de todo o Brasil foi 52,5.
A menor Média nessa questão foi obtida pelos alunos da região Nordeste ( 44,8 ), enquanto a maior Média foi obtida na região Sul ( 59,3 ).
Quanto à variabilidade das notas, o Desvio padrão de todo o Brasil foi 27,7.
O maior Desvio padrão foi obtido na região Sul ( 27,6 ), enquanto o menor foi obtido na região Nordeste ( 26,5 ).
A nota Máxima, 100,0 pontos, foi alcançada por, pelo menos, um aluno de todas as regiões, a menos da região Nordeste na qual a nota Máxima foi 90,0.
A Mediana do Brasil como um todo foi 55,0, a mesma para as regiões Sudeste e Centro-Oeste.
A nota Mínima ( 0,0 ) foi obtida por, pelo menos, um estudante em todas as regiões do Brasil.
O Gráfico 3.11 mostra a distribuição das notas na questão discursiva 3, do Componente de Conhecimento Específico da área de Tecnologia em Agronegócio.
O intervalo de maior frequência foi o ( 60, 70 ].
O intervalo ( 40, 50 ] aparece como um máximo local, assim como o grupo daqueles que obtiveram zero, com cerca de 8 \% do total de respostas. 3.3.2.2 - Comentários sobre as respostas à Questão Discursiva 3 A questão aborda um dos temas principais da compreensão da estrutura do agronegócio e seu funcionamento, e se relaciona às Diretrizes Curriculares Nacionais dos cursos de tecnólogos em agronegócios.
O tema é bastante complexo por envolver conhecimento amplo de todas atividades relativas ao setor agropecuário, desde a produção de insumos, a produção agropecuária, a colheita, o armazenamento, o processamento, até o setor de serviços relacionados.
Sem dúvida, os concluintes dos cursos devem ter a compreensão do fluxo de uma cadeia produtiva para promoção de melhorias na sua eficiência, gerando melhor custo/benefício para o produto final.
A questão apresentou dois itens.
No item ‘ a ’ o estudante deveria escolher uma cadeia produtiva e explanar sobre as suas diferentes etapas e atividades, a fim de o estudante demonstrar pleno conhecimento de como é formada uma cadeia produtiva no agronegócio e as suas inter-relações, também chamadas de ‘ elos ’ da cadeia produtiva.
No item ‘ b ’, o estudante deveria explicar como uma determinada cadeia produtiva ( do milho ) poderia influenciar em outras cadeias no agronegócio.
Nesse item, era necessário ao estudante explorar um conhecimento mais amplo de como ocorrem intercomunicações entre diferentes cadeias que se relacionam, influenciando no abastecimento, no mercado, nos preços, entre outros fatores relacionados ao agronegócio.
A questão era bem direta e clara, e apontava para respostas igualmente diretas.
De certa forma, permitia que o estudante pudesse se expressar de maneira personalizada sobre o assunto, mostrando seu conhecimento pessoal sobre o tema.
O enunciado não dava margem a interpretações diferentes da esperada pelo padrão de respostas, a não ser que o estudante não soubesse responder por falta de conhecimentos sobre a área abordada na questão.
Merece destaque que as respostas ao item ‘ b ’ possibilitavam muitas opções de correlação entre cadeias produtivas, o que foi considerado no padrão de respostas.
Observou -se que os estudantes utilizaram quase todo o espaço para suas respostas, principalmente, nas mais completas, mas isso também ocorreu com as respostas erradas.
Além de o assunto ser bastante amplo e haver diferentes formas de abordagem, a caligrafia de alguns estudantes não permitiu respostas mais longas e/ou elaboradas no pequeno espaço disponível na prova.
A questão pode ser considerada de dificuldade média, em comparação com as demais questões de formação específica, já que é um tema bastante comum e bem abordado nos currículos dos cursos em agronegócios.
No entanto, exigia que o aluno tivesse conhecimento da complexidade que uma cadeia produtiva pode atingir, influenciando outras atividades e setores a ela relacionados.
No item ‘ a ’, houve uma maior quantidade de acertos de acordo com o padrão de resposta proposto.
Em alguns casos, houve sobreposição de respostas entre o item ‘ a ’ e o item ‘ b ’, mais devido à falta de conhecimento sobre o item ‘ a ’ do que de buscar algum exemplo diferente de cadeia produtiva.
Houve exemplificação de diferentes cadeias produtivas, no entanto, nem sempre, o estudante conseguiu relacionar os três grandes ‘ elos ’ : antes da porteira, dentro da porteira e depois da porteira, de forma correta e clara.
No item ‘ b ’, a maioria dos estudantes respondeu diretamente que a cadeia produtiva de milho interferia na produção de animais, não sabendo explicar de que forma.
Poucos conseguiram explicar, de forma mais clara, que a cadeia do milho afetaria a produção de insumos básicos na produção de rações para animais, interferindo, indiretamente, nos preços da carne e do leite, por exemplo, nas entressafras e na produção de alimentos para consumo humano como óleos vegetais, farinhas e derivados.
Houve, algumas vezes, confusão entre termos específicos da área, trocas de conceitos e definições acerca do tema abordado.
Apesar disso, na maior parte das \textbf{repostas}, a linguagem adotada pelos alunos foi adequada no uso dos termos específicos da área.
Observaram -se, também, casos de alunos que apenas responderam a um dos dois itens, indicando pressa na leitura do enunciado e em responder à questão.
Os estudantes apresentaram um desempenho mediano nessa questão.
Demonstraram bastante dificuldade de escrever um texto sucinto e expor todo o tema abordado na questão.
Dificuldade na escrita e caligrafia ruim foram desafios para a boa compreensão das respostas.
Os que não utilizaram todo o espaço disponível para responder à questão o fizeram, ora pela falta de uma resposta mais completa ora pela falta de resposta literalmente, já que houve muitos que não responderam a todos os itens.
Respostas fracas, aproximadamente 15 \% do total, foram aquelas que tratavam de apenas um item e de forma incompleta.
Respostas medianas representaram mais de 50 \% do total e, nelas, apesar de não abrangeram todos os detalhes dos itens solicitados, o que contemplaram foi feito de forma satisfatória.
Respostas boas, aproximadamente 25 \% do total, foram as mais completas e de acordo com o padrão.
De modo geral, ficou demonstrada uma grande dificuldade na interpretação de conceitos abordados na questão.
Houve muita confusão entre conceitos de cadeia produtiva e de sistemas de produção agropecuária, indicando deficiência dos estudantes na compreensão do assunto avaliado.
Ocorreram, também, estudantes que buscaram a junção das respostas aos itens ‘ a ’ e ‘ b ’, na tentativa de aproveitar a mesma cadeia produtiva para explicar o que foi solicitado em cada item.
Percebeu -se, em determinadas respostas, alguma deficiência de conhecimento sobre o assunto, no entanto, não foi possível identificar se esta deficiência advém da falta de conhecimento ou, simplesmente, da dificuldade em responder a uma questão discursiva.
Em algumas respostas, não se demonstrava preocupação em dar sequência ao raciocínio ou fazer a relação entre os temas que foram solicitados, faltando um maior comprometimento no encadeamento das ideias.
Verificaram -se quantitativos inferiores a 10 \% de questões em branco e a 1 \% de questões desconsideradas. 3.3.2.3 - Análise da Questão Discursiva 4 do Componente de Conhecimento Específico A Tabela 3.19 contém as informações relativas à questão 4 do conjunto de questões do Componente de Conhecimento Específico.
O desempenho dos estudantes nessa questão foi superior ao das questões 3 e 5.
A Média geral do Brasil foi 53,5, sendo a menor Média registrada na região Norte ( 51,0 ), e a maior, na região Nordeste ( 55,8 ).
A nota Máxima ( 100,0 ) foi atingida por, pelo menos, um concluinte de todas as regiões, com exceção da região Norte que teve note Máxima 80,0.
A Mediana do Brasil foi 60,0, o mesmo valor obtido em todas as regiões.
A nota Mínima do Brasil foi zero e foi encontrada em todas as regiões.
O Gráfico 3.12 representa a distribuição de notas da questão discursiva 4, do Componente de Conhecimento Específico.
Essa distribuição tem moda no intervalo ( 50 ; 60 ], com mais da metade das respostas.
O próximo intervalo de maior frequência abarca menos do que 20 \% das respostas. 3.3.2.4 - Comentários sobre as respostas à Questão Discursiva 4 A questão se relaciona às Diretrizes Curriculares Nacionais dos cursos de tecnólogos em agronegócios e aborda tema relativo à melhoria da qualidade de vida do trabalhador rural.
Este tema é obrigatório nos currículos de tecnólogos e engenheiros ligados ao setor agropecuário.
A temática trouxe à tona assuntos relacionados a campanhas educativas e de melhoria do ambiente de trabalho, além da legislação pertinente.
O enunciado era claro e adequado para interpretação pelo estudante, que deveria mostrar conhecimento geral e básico sobre o assunto envolvido.
A questão possibilitava uma gama bastante ampla e diversificada de respostas, dentro do tema proposto.
Por essa razão, a questão oferecia, também, boa possibilidade de avaliação do \textbf{aprendizado}.
Com respeito à legislação, deve -se considerar que, tecnólogos, provavelmente, devem ter um conhecimento bem básico deste assunto, já que vão atuar de forma conjunta com engenheiros e/ou outros profissionais com melhor qualificação e maior atribuição profissional para tal atividade.
A questão pode ser considerada como de dificuldade mediana, considerando que propõe ao aluno abordar, de forma dissertativa, o tema de prevenção de riscos ocupacionais no ambiente de trabalho rural e mencionar a legislação pertinente.
No entanto, cabe chamar atenção para o fato de que, nem sempre, é dada grande importância e valor a esse tema dentre os mais diferentes conteúdos que fazem parte das ementas das disciplinas do curso.
Dessa forma, ocorreu, sempre, uma certa dificuldade dos alunos em responder, de forma adequada, ao solicitado pela questão.
As respostas dadas pelos estudantes, na sua maioria, não foram completas e abordaram, apenas, o uso de EPI e/ou alguma outra técnica isoladamente.
Foram poucas as respostas mais completas que abordaram todos os aspectos contemplados no padrão de respostas.
O assunto não foi explicado ou abordado, de forma adequada, tanto que, muitas vezes, não foi utilizado todo o espaço destinado à resposta.
Novamente, observa -se pouca preocupação do estudante em escrever de forma correta, ter boa caligrafia e desenvolver, de forma concisa, as explicações para a resposta da questão proposta.
O assunto parece não ter sido assimilado ou ter tido pouca ênfase no currículo dos estudantes que fizeram a prova, pelo que pode ser depreendido da maioria das respostas observadas.
Poucos alunos conseguiram associar e comentar suas respostas segundo a legislação pertinente, demonstrando pouco ou superficial conhecimento na área do tema proposto pela questão.
Os estudantes apresentaram um desempenho mediano nessa questão.
Observou -se grande dificuldade em responder à questão abordando todos os itens do padrão de resposta.
Houve grande quantidade de respostas em branco, o que talvez resulte de pouco estudo acerca do assunto abordado.
A maioria das respostas eram sucintas, sem demonstrar conhecimento das relações entre trabalhador e empregador na atividade rural.
Respostas fracas, que representaram menos de 5 \% do total, foram aquelas em que o aluno não respondeu a nenhuma ação preventiva ou não entendeu o que foi perguntado.
Respostas medianas, que representaram mais de 70 \% do total, foram aquelas em que os alunos responderam, de forma satisfatória, com pelo menos uma ação preventiva, explicando sua vinculação com o risco ocupacional.
Respostas boas, que representaram, aproximadamente, 15 \%. aquelas em que os alunos responderam, de forma satisfatória, com pelo menos uma ação preventiva, explicando sua vinculação com o risco ocupacional.
Respostas boas, que representaram, aproximadamente, 15 \%.
De um modo geral, observou -se que a maioria das respostas abordava o uso de EPI como a principal forma de diminuir riscos de acidente de trabalho.
Em poucos casos, o estudante discutiu o uso de outras técnicas e ações para a redução de acidentes de trabalho rural e, muito menos esclareceu algumas ações apontadas para promover melhorias nas condições de trabalho no campo.
Observou -se falta de conhecimento do estudante sobre ações para o empregador minimizar os riscos ocupacionais, de forma preventiva, no ambiente de trabalho.
Foi grande o número de respostas que indicaram apenas uma das formas de minimizar riscos no ambiente de trabalho rural.
Provavelmente, isso se deve à falta de conhecimento sobre a legislação pertinente ou à dificuldade de visualizar diferentes tipos de riscos no ambiente de trabalho rural.
As respostas dadas evidenciaram dificuldade de organização de ideias e de elaboração de textos que demonstrassem domínio do conhecimento requerido na questão. 3.3.2.5 - Análise da Questão Discursiva 5 do Componente de Conhecimento Específico A Tabela 3.20 contém as informações relativas à questão 5 do conjunto do Componente de Conhecimento Específico.
O desempenho dos estudantes de todo o Brasil nesta questão foi o mais baixo dentre as três questões discursivas desse componente.
A nota Média dos estudantes de todo o Brasil foi 26,4.
A maior Média foi registrada na região Sul ( 31,4 ), enquanto a menor Média foi registrada na região Nordeste ( 22,3 ).
Quanto à variabilidade das notas, o Desvio padrão dos alunos do Brasil, como um todo, foi 26,5.
Enquanto o maior desvio foi encontrado na região Sudeste ( 27,5 ), o menor foi encontrado na região Nordeste ( 23,9 ).
A Mediana para o Brasil foi 20,0, a mesma de quatro das cinco regiões ( a exceção foi a região Sul que teve Mediana 30,0 ) Para o conjunto de alunos de Tecnologia em Agronegócio do Brasil, a nota Máxima foi 100,0, e a nota Mínima foi zero.
Os mesmos limites foram encontrados em todas as regiões O Gráfico 3.13 representa a distribuição das notas da questão discursiva 5 do Componente de Conhecimento Específico.
Essa distribuição tem moda no intervalo ( 10 ; 20 ], seguido de perto pelo grupo daqueles que obtiveram zero.
Destaca -se também, o grande número de estudantes que deixaram esta questão em branco.
A soma da frequência desse grupo com a frequência daqueles com zero perfaz mais do que 35 \% do total de respostas. 3.3.2.6 - Comentários sobre as respostas à Questão Discursiva 5 O assunto abordado na questão era atual e relevante para o setor agropecuário nacional.
A sustentabilidade dos sistemas produtivos é uma questão chave para o desenvolvimento do país e importante para a manutenção da capacidade produtiva no campo.
A questão traz o debate sobre a conservação da natureza e a manutenção da produção agrícola.
Por sua importância e atualidade é um assunto que faz parte dos currículos nacionais de cursos da área de ciências agrárias e deve ser visto como prioridade nas atividades do egresso do curso de agronegócios.
O estudante deveria mostrar conhecimento básico do assunto e descrever a importância do Cadastro Ambiental Rural ( CAR ) para a conservação dos recursos naturais na propriedade agrícola, correlacionando -o com a melhoria e manutenção da produtividade rural.
De forma geral, como o assunto ainda é recente, há posições contraditórias em discussão no meio acadêmico e no meio rural.
Dessa forma, verificaram -se, por exemplo, respostas favoráveis e não favoráveis à implantação do CAR.
No entanto, seria necessário que o estudante reconhecesse que a proposta do CAR possui um impacto positivo sobre o sistema de produção rural, uma vez que o produtor deve adequar sua propriedade ambientalmente, por meio de uma declaração, o que possibilita uma fiscalização mais efetiva.
O principal objetivo do cadastro é gerar uma maior consciência ambiental para os proprietários e trabalhadores na atividade rural do país.
A questão pode ser considerada de dificuldade mediana, já que busca avaliar conhecimento básico e atualizado de assunto de grande importância no contexto do agronegócio brasileiro e da sustentabilidade da atividade.
A finalidade do CAR e seus impactos no agronegócio são definições que devem estar claras e sedimentadas para o profissional que vai atuar na atividade rural do país.
As respostas, em sua maioria, não abordavam o tema de forma abrangente.
Houve grande quantidade de respostas em branco, e muitas com pouco texto redigido, sem desenvolver o assunto de forma adequada.
Aparentemente, há grande deficiência na compreensão dos benefícios que o cadastro traz para o produtor rural e para o país como um todo.
Geralmente, o estudante descrevia uma das finalidades do CAR, mas não explicava seus impactos para o agronegócio, desconsiderando ou desconhecendo sua importância para : a manutenção dos recursos naturais, a verificação da legislação ambiental, a gestão ambiental da propriedade, a fiscalização, a integração de A maior parte das respostas foi superficial e com pouca explanação sobre as funções e finalidades do CAR.
Como é um tema regulamentado e implantado há pouco tempo, ainda existe grande resistência, por parte da comunidade rural, para reconhecer os benefícios da adoção do CAR e, talvez, por isso, ao abordá -lo nos currículos de ensino, a ênfase recaia na polêmica.
O CAR é, também, uma ferramenta por meio da qual o proprietário passa a se autofiscalizar e gerenciar melhor a sua atividade rural, de forma a respeitar a legislação, auxiliar a fiscalização e planejar de forma sustentável novas atividades agrícolas.
Os estudantes apresentaram um desempenho fraco nessa questão.
Muitos textos estavam mal escritos, com caligrafia ruim, com respostas incompletas e superficiais.
Houve grande quantidade de respostas em branco e, mesmo nas \textbf{repostas} mais completas, raramente se percebia reconhecimento da importância do CAR para a atividade rural.
Respostas fracas, que representaram mais de 50 \% do total, foram aquelas em que o aluno não conseguiu apontar, de alguma forma, uma finalidade do CAR, explicando a integração de informações e/ou um ou mais impactos do cadastro para o agronegócio.
Respostas medianas, que representaram mais de 25 \% do total, foram aquelas em que o aluno apenas indicou algum impacto e/ou finalidade, no entanto, sem explicar como ocorre esta integração.
Respostas boas, que representaram pouco mais de 5 \% do total, foram aquelas que conseguiram descrever, de forma mais dissertativa e completa, o que foi solicitado no enunciado da questão.
Como já apontado anteriormente, de forma geral, as respostas dos estudantes foram evasivas e não comentaram as diferentes funções do CAR e seus principais impactos para o agronegócio brasileiro.
Verificou -se que alguns estudantes desconhecem o assunto, principalmente pelo número de respostas em branco e ênfase nos malefícios que poderiam decorrer da adoção do cadastramento.
Muitos estudantes enfatizaram a perda de produção agrícola na propriedade pela manutenção das áreas de preservação, que são obrigatórias pela legislação, antes mesmo da institucionalização do CAR, indicando, equivocadamente, que o cadastro não traz benefício ao agronegócio.
A maioria das respostas era curta, sem demonstrar embasamento teórico.
Mais uma vez ficou clara a grande dificuldade de elaboração de respostas dissertativas.
O grande número de respostas em branco ( em torno de 10 \% ) pode indicar que o assunto não é conhecido ou que talvez não tenha sido abordado de forma adequada em sala de aula.
Por outro lado, as respostas dadas pelos estudantes também revelam como o tema é tratado dentro e fora do meio acadêmico.
Foram corrigidas respostas indicando impactos benéficos e malefícios na adoção do CAR pelas propriedades rurais.
A implantação recente do CAR e as mudanças na legislação pertinente podem ser fatores intervenientes no grau de esclarecimento, no tipo de abordagem feita no meio acadêmico e, consequentemente, no processo que conduz à demonstração da importância que o assunto possui para a atividade rural do país. 3.3.3 - Considerações Finais De um modo geral, verificou -se um desempenho mediano dos alunos nas questões discursivas específicas da área.
Existem pontos que merecem ser observados como a dificuldade dos alunos para redigir respostas para questões discursivas, abordando temas amplos de forma concisa e, sobretudo, demonstrando conhecimento sobre os conteúdos avaliados pelas questões.

% matched lemmas: aprendizado, colonização, repor
\end{textsample}
